\par\Needspace{9\baselineskip}\bigskip
\pdfbookmark[3]{Darstellung, [25]}{pallas-II-article25}
\begin{center}[25.]\end{center}
\begingroup
\noindent\hspace*{\parindent}Wir geben nun eine sehr einfache Methode an, die Störungen von der Ordnung des Quadrats der störenden Kraft vollständig zu berechnen. Es sei \(Q\mathrm{d}t\) die momentane Änderung eines Elements, zum Beispiel von \(\AscendingNode+\delta\AscendingNode\). Dann hängt \(Q\) von den Elementen des gestörten und des störenden Planeten ab. Setzt man die in erster Ordnung der Störungen richtigen Werte dieser Elemente ein, so ist klar, dass \(Q\) bis zur zweiten Ordnung (einschließlich) richtig ist. Es seien also \SourceLeaf{475}\(a,\ b,\ c,\ e,\ f,\ g\) die mittleren Werte der Elemente des gestörten Planeten und \(a+\Delta a,\ b+\Delta b,\ c+\Delta c,\ e+\Delta e,\ f+\Delta f,\ g+\Delta g\) die in erster Ordnung richtigen Werte. Ferner mögen \(a',\ b',\ c',\ e',\ f',\ g',\ a'+\Delta a'\) usw. dieselben Größen für den störenden Planeten bezeichnen. Man kann ansetzen
\[
\begin{alignedat}{2}
Q=Q^0&{}+\Delta a\cdot\left(\frac{\mathrm{d}Q^0}{\mathrm{d}a}\right)&&{}+\Delta a'\cdot\left(\frac{\mathrm{d}Q^0}{\mathrm{d}a'}\right)\\[3pt]
&{}+\Delta b\cdot\left(\frac{\mathrm{d}Q^0}{\mathrm{d}b}\right)&&{}+\Delta b'\cdot\left(\frac{\mathrm{d}Q^0}{\mathrm{d}b'}\right)\\[3pt]
&{}+\Delta c\cdot\left(\frac{\mathrm{d}Q^0}{\mathrm{d}c}\right)&&{}+\Delta c'\cdot\left(\frac{\mathrm{d}Q^0}{\mathrm{d}c'}\right)\\[3pt]
&\quad\text{usw.}&&\quad\text{usw.,}
\end{alignedat}
\]
\noindent wenn man sich auf die Größen zweiter Ordnung beschränkt. Die Größen \(\Delta a,\allowbreak\ \Delta b,\allowbreak\ \Delta c,\allowbreak\ \Delta e,\allowbreak\ \Delta f,\allowbreak\ \Delta g\) enthalten jeweils ein der Zeit proportionales Glied und periodische Funktionen von \(nt\) und \(n't\) oder .\ .\ .\ .\ .
\par\medskip\begin{center}\rule{.18\linewidth}{.3pt}\end{center}
\endgroup
